\section{Remember Thy Last End}

November begins among the dead. On its first day the Church keeps the feast of all the saints, and on its second she remembers all the faithful departed; the Office of the Dead is said, and at its first lesson the Church takes the voice of Job and prays, \latin{Parce mihi, Domine: nihil enim sunt dies mei}, ``spare me, for my days are nothing.''\endnote{Job 7:16, Douay--Rheims (Challoner). Latin: \emph{Rituale Romanum} (Ratisbon: Pustet, 1872), Officium Defunctorum, Matins, first nocturn, Lectio I.} The year is running out. On the Last Sunday after Pentecost, in the 1962 Missal, the Gospel of the Mass turns to the end of all things, to the abomination of desolation and the sign of the Son of man in heaven, and closes on the promise, ``Heaven and earth shall pass: but my words shall not pass.''\endnote{Matthew 24:15--35, the Gospel of the Last Sunday after Pentecost in the 1962 Missal; the sentence quoted is 24:35, Douay--Rheims (Challoner).} It is the season in which the Church sets before her children the four last things: death, judgment, hell, and heaven. The first of them is death, because it is the door to the other three, and because it is the one that every man will certainly pass.

The counsel that gave these things their name stands near the end of a chapter of practical wisdom in the book of Ecclesiasticus, and it is worth reading in the company the sage gave it. ``A gift hath grace in the sight of all the living, and restrain not grace from the dead. Be not wanting in comforting them that weep, and walk with them that mourn. Be not slow to visit the sick: for by these things thou shalt be confirmed in love. In all thy works remember thy last end, and thou shalt never sin.''\endnote{Ecclesiasticus 7:37--40, Douay--Rheims (Challoner).} The three verses before it are works of mercy: kindness to the dead, comfort to the mourners, the visit to the sick. The fourth names what such works teach. The Latin says \latin{memorare novissima tua}, ``remember thy last things,'' and from that plural the Church's teachers drew the name \latin{novissima}, the last things, for the realities that await every soul at the end of its course.\endnote{Ecclesiasticus 7:40 in the Clementine Vulgate: \latin{In omnibus operibus tuis memorare novissima tua, et in \ae ternum non peccabis}.} So the remembrance of death is commended here at a bedside and a graveside. It is learned by keeping company with the dying and the bereaved, and it bears its fruit in a life that does not sin.

The Church writes the same counsel on the forehead. On Ash Wednesday, at the threshold of Lent, the priest puts blessed ashes on the heads of the clergy and people and says to each one, \latin{Memento, homo, quia pulvis es, et in pulverem reverteris}: ``Remember, O man, that thou art dust and unto dust thou shalt return.''\endnote{The formula at the imposition of ashes on Ash Wednesday, in F. X. Lasance, \emph{The New Roman Missal} (New York: Benziger, revised edition of 1945), Latin and English.} The words are God's own, spoken to Adam as he went out of paradise: ``In the sweat of thy face shalt thou eat bread till thou return to the earth out of which thou wast taken: for dust thou art, and into dust thou shalt return.''\endnote{Genesis 3:19, Douay--Rheims (Challoner).} What the liturgy says once a year in Lent, the saints learned to say every day. Alphonsus Liguori opens his \emph{Preparation for Death} on this verse, and his first consideration begins with it plainly: ``Consider that you are dust, and that you shall return to dust. A day will come when you shall die, and rot in a grave.''\endnote{Alphonsus Liguori, \emph{Preparation for Death}, Consideration I, first point, in the English of Dublin: James Duffy, 1843.}

The promise attached to this remembrance is a large one: \emph{thou shalt never sin}. The Fathers and the masters of the spiritual life did not take it for a figure of speech. They held that a man who truly keeps his death before him will order his life by it, and over the centuries they made the remembrance a daily discipline, a meditation for retreats, and at last a monthly exercise. The masters of the early modern centuries, Francis de Sales and Alphonsus among them, wrote for people in the world as well as for monks, and they taught the meditation in a fixed order: first the truth of death, as faith knows it; then its consideration, until the truth is felt; then the resolution it calls for, and the prayer in which the resolution is kept. The Church follows that order in her own year. She remembers the dead in November, prays the Requiem, watches through the last Sundays for the coming of her Lord, and opens Advent with the command to watch.

Death, then, is not proposed to the Christian as a riddle for the intellect or as a terror to be endured in silence. It is proposed as something to be remembered, prepared for, and in the end received from the hand of God. The Church has a great deal to say about it, because she has kept watch at deathbeds since the days of the Apostles, and her words are words of faith, of fear, and of a hope that does not deceive.
